\documentclass[10pt,a4paper]{amsart}
\usepackage[margin=19mm]{geometry}
\usepackage{amsmath,amssymb,mathtools}
\usepackage[scr=rsfs,cal=euler]{mathalfa}
\usepackage{xfrac}
\usepackage[unicode,hidelinks,bookmarks=false]{hyperref}
\newcommand{\ie}{i.e.\ }
\newcommand{\cf}{cf.\ }
\newcommand{\resp}{respectively\ }
\newcommand{\Subsection}{\subsection}
\providecommand{\nobreakdash}{\nobreak\hbox{-}}
\setlength{\emergencystretch}{2em}
\multlinegap0pt
\usepackage{mathtools}
\usepackage{tikz}
\newtheorem{lemma}{Lemma}
\DeclareMathOperator{\Prob}{Pr}
\newcommand{\mbf}[1]{\mathbf{#1}}
\title{Constructive Tchakaloff results and well-conditioned quadrature through randomized least squares}
\author{Filip B\v{e}l\'{i}k, Akil Narayan, John Turnage}
\date{Source: arXiv:2609.09506v1 (CC BY 4.0)}
\begin{document}
\maketitle

\noindent\emph{Source and licence.} Constructive Tchakaloff results and well-conditioned quadrature through randomized least squares. Version arXiv:2609.09506v1 (CC BY 4.0), distributed by arXiv under the Creative Commons Attribution 4.0 International licence, http://creativecommons.org/licenses/by/4.0/, as recorded in the arXiv licence metadata for this exact version and shown on the abstract page https://arxiv.org/abs/2609.09506v1. Full licence notice: \texttt{licenses/arxiv-2609.09506-CC-BY-4.0.txt}.\par\smallskip

\noindent\textit{Editorial scope.} Counted unit: the article's lemma lemma:admissible-M-window (source0.tex lines 1883--1928). The article's complete original proof of it is delivered with it (source0.tex lines 1930--1997, one \texttt{proof} environment of 68 lines). No mathematics is written, repaired or re-derived: the statement and the proof are the article's own lines, in the article's own notation, and no retained line is substituted or re-flowed.  Retained uncounted as the source-local context the proof needs: lines 1610--1627.  Nothing was omitted from any retained span, so the package carries no omission record.  No retained line contains a citation, so the wrapper carries no bibliography and no bibliography entry.  No retained line contains a figure environment, an image-inclusion command or drawing code, so no image is delivered and the package declares no asset dependency.  Editorial formatting only, in addition: an AMS \texttt{amsart} preamble that restates the article's own declarations and macros used by the retained lines, namely the environments \texttt{lemma} on separate counters, together with the standard packages those lines need.  The retained environments print in this document as Lemma 1 in order of appearance, whereas the article prints the same items with the numbering of its own full document.  The complete native source of the article as distributed in the arXiv e-print for this exact version is bundled unchanged as \texttt{sources/209819/source0.tex}.
\begin{lemma}[Inverse Gramian concentration]
\label{lemma:F-chernoff}
    Let \(c_0 \coloneqq 1 - 2^{-1/2}\). Suppose that \(Q < \infty\), \(\delta\in(0,c_0]\), and \(\varepsilon>0\). If
    \begin{align}
        \label{eq:Fhat-sampling-criterion}
        M
        \ge
        \frac{6Q}{\delta^2}
        \log\left(\frac{2N}{\varepsilon}\right)
        \implies
        \Prob\left[
            \left\|
            \left(\mbf G^{(M)}\right)^{-1}-\mbf I
            \right\|_2>\delta
        \right]
        \leq \varepsilon .
    \end{align}
\end{lemma}

\begin{lemma}[Admissible sample-size window]
\label{lemma:admissible-M-window}
Let \(0<\varepsilon,\beta<1\), \(\sigma>0\), and \(J\geq1\). Suppose that
\(r_J\leq1/2\),
\begin{equation}\label{eq:window-compatibility-condition}
    \frac{QJ}{\sigma^2}
    \log\left(\frac{2N}{\varepsilon}\right)
    \geq1,
    \qquad 
    \text{and}
    \qquad
    \frac{14QJr_J}{\sigma^2}
    \log\left(\frac{2N}{\varepsilon}\right)
    <
    \log\left(\frac{1}{1-\beta}\right).
\end{equation}
Then there exists an integer \(M\) satisfying
\begin{equation}\label{eq:window-length-condition}
    \frac{6QJ}{\sigma^2}
    \log\left(\frac{2N}{\varepsilon}\right)
    \leq M
    \leq
    \frac{7QJ}{\sigma^2}
    \log\left(\frac{2N}{\varepsilon}\right)
\end{equation}
such that
\(
    \Prob[R_M(J)]>1-\beta.
\)
If, in addition,
\(
    \frac{\sigma}{\sqrt J}\leq c_0,
\)
then
\[
    \Prob\left[
        \left\|
            \left(\mbf G^{(M)}\right)^{-1}-\mbf I
        \right\|_2
        \leq
        \frac{\sigma}{\sqrt J}
    \right]
    \geq
    1-\varepsilon.
\]
\end{lemma}

\begin{proof}
Set
\[
    B
    \coloneqq
    \frac{QJ}{\sigma^2}
    \log\left(\frac{2N}{\varepsilon}\right).
\]
By assumption, \(B\geq1\), so the interval \([6B,7B]\) contains an integer \(M\). 
Now, since the samples are independent,
\(
    \Prob[R_M(J)]
    =
    (1-r_J)^M.
\)
Moreover, \(r_J\leq1/2\) implies
\(
    -\log(1-r_J)\leq2r_J.
\)
Therefore,
\begin{align*}
    -\log\Prob[R_M(J)]
    &=
    -M\log(1-r_J) 
    \leq
    2Mr_J 
    \leq
    14Br_J 
    <
    \log\left(\frac{1}{1-\beta}\right).
\end{align*}
Exponentiating gives
\(
    \Prob[R_M(J)]>1-\beta.
\)
Finally, suppose that
\(
    \frac{\sigma}{\sqrt J}\leq c_0
\)
and set
\(
    \delta
    =
    \frac{\sigma}{\sqrt J}.
\)
The lower bound on \(M\) gives
\[
    M
    \geq
    \frac{6QJ}{\sigma^2}
    \log\left(\frac{2N}{\varepsilon}\right)
    =
    \frac{6Q}{\delta^2}
    \log\left(\frac{2N}{\varepsilon}\right).
\]
Hence, Lemma \ref{lemma:F-chernoff} yields
\[
    \Prob\left[
        \left\|
            \left(\mbf G^{(M)}\right)^{-1}-\mbf I
        \right\|_2
        \leq
        \frac{\sigma}{\sqrt J}
    \right]
    \geq
    1-\varepsilon.
\]
\end{proof}

\end{document}
